\chapter{Abelian Groups}
\section{Classification of Finitely Generated Abelian Groups}

We abbreviate the following theorem by \textbf{FTFGAG}. 
\begin{theorem}[Fundamental Theorem of Finitely Generated Abelian Groups]
  Let \(G\) be a finitely generated abelian group. Then \begin{itemize}
    \item there exist primes \(p_1 < \cdots < p_k\) and decreasing positive integers \(\{n_{ij}\}_j\) such that \[
      G \cong \Z^r \oplus \bigoplus_{i,j}\Z_{p_i^{\smash[b]{n_{ij}}}};
    \]
    \item there exist positive integers \(1 < d_1 \mid d_2 \mid \ldots \mid d_s\) such that \[
      G \cong \Z^r \oplus \Z_{d_1} \oplus \Z_{d_2} \oplus \ldots \oplus \Z_{d_s}.
    \]
    Furthermore, these two decompositions are uniquely determined.
  \end{itemize}
\end{theorem}

The integer \(r\) is called the \emph{rank} of \(G\). The primes \(p_i^{n_{ij}}\) are called the \emph{elementary divisors} of \(G\). The integers \(d_i\) are called the \emph{invariant factors} of \(G\).

The structure theorem applied to finite groups sets \(r = 0\), which gives a bijective correspondence \begin{align*}
  \left\{
    \begin{aligned}
        &\text{partitions of \(n_i\)} \\
        &\text{for \(i = 1,\ldots,k\)}
    \end{aligned}\right\}
    \xleftrightarrow{\hspace*{5em}}
  \left\{
    \begin{aligned}
        &\text{finite abelian groups} \\
        &\text{of order \(p_1^{n_1}\cdots p_k^{n_k}\)} \\
  \end{aligned}\right\}
\end{align*}

\begin{corollary}
  Let \(\K\) be a field and put \(\K^\times = \K\setminus\{0\}\).
  Every finite group \(G < \K^\times\) is cyclic.
\end{corollary}
\begin{proof}
  By the structure theorem of finitely generated abelian groups, we decompose \[
    G \cong \Z_{d_1} \oplus \Z_{d_2} \oplus \ldots \oplus \Z_{d_s},
  \] where \(d_i = p_i^{n_i}\) for some prime \(p_i\) and positive integer \(n_i\).
  Let \(m = \lcm(d_1,\ldots,d_s) \leq |G|\), then \(a^m = 1\) for all \(a \in G\).
  However, the polynomial \(x^m - 1\) has at most \(m\) roots in \(\K\), so \(|G|\leq m\).
  Since \[
    \lcm(d_1,\ldots,d_s) = |G| =  d_1\cdots d_s,
  \] we have \(d_1,\ldots,d_s\) pairwise coprime, then by the Chinese remainder theorem, we have \[
    G \cong \Z_{d_1}\oplus\cdots\oplus\Z_{d_s} \cong \Z_{d_1\cdots d_s} \cong \Z_{|G|}.
  \]
\end{proof}

\begin{example}
  \(\Z_{13}^\times = \{1,2\ldots,12\} = \<2\> = \Z_{12}\).
\end{example}
\begin{proof}
  Since \(13\) is prime, every nonzero residue modulo \(13\) is invertible, so \(\Z_{13}^\times = \{1,2,\ldots,12\}\) under multiplication modulo \(13\).

  The successive powers of \(2\) modulo \(13\) are:
  \[
    \begin{array}{c|ccccccccccccc}
    k&0&1&2&3&4&5&6&7&8&9&10&11&12\\
    \hline
    2^k\bmod 13&1&2&4&8&3&6&12&11&9&5&10&7&1
    \end{array}
  \]
  Hence \(\Z_{13}^\times = \<2\> \cong \Z_{12}\).
\end{proof}

\begin{example}
  \(\Q\) is not finitely generated as an abelian group.
\end{example}
\begin{proof}
  Suppose \(Q = \<q_1,\ldots,q_n\>\), then we write \(q_i = a_i/b_i\). Put \(B = b_1\cdots b_n\), then since \(1/2B\in\Q\), we can decompose \[
    \frac{1}{2B} = k_1\frac{a_1}{b_1} + \cdots + k_n\frac{a_n}{b_n}.
  \]
  Then we obtain \[
    \frac{1}{2} = k_1a_1\frac{B}{b_1} + \cdots + k_na_n\frac{B}{b_n} \in \Z,
  \] a contradiction.
\end{proof}

\begin{proposition}
  Let \(G\) be group and \(N\) a normal subgroup of \(G\). Then \begin{enumerate}
    \item if \(G\) is cyclic, then \(G/N\) is cyclic;
    \item if \(G\) is abelian, then \(G/N\) is abelian;
  \end{enumerate}
\end{proposition}
\begin{proof}
  \begin{enumerate}[label=(\alph*)]
    \item Suppose \(G = \<a\>\) for some \(a\in G\setminus\{e\}\). Then we show \(G/N = \<aN\>\). Let \(bN\in G/N\), then \(b = a^k\) for some \(k\), then \[
      bN = a^kN = (aN)^k \in \<aN\>.
    \]
    \item Suppose \(G\) is abelian. Let \(aN, bN\in G/N\), then \[
      (aN)(bN) = (ab)N = (ba)N = (bN)(aN). \qedhere
    \]
  \end{enumerate}
\end{proof}

\begin{example}
  Consider \(G = \Z_4\times\Z_6\) and \(N = \<(2,3)\>\lhd \Z_4\times\Z_6\).
  We compute first \[
    [G:N] = \frac{|G|}{|N|} = \frac{4\cdot 6}{2} = 12.
  \]
  By the structure theorem, there are only 2 abelian groups of order \(12\) up to isomorphism, namely \[
    \Z_2\times\Z_6 \quad\text{and}\quad \Z_{12}.
  \]
  Consider the coset \((1,0) + N\), since \(N = \{(0,0), (2,3)\}\), we see \[
    2(1,0) = (2,0) \notin N, \quad
    3(1,0) = (3,0) \notin N, \quad
    4(1,0) = (0,0) \in N,
  \] so \(|(1,0) + N| = 4\).

  Since in \(\Z_2\times\Z_6\), \[
    |(a,b)| = \lcm(|a|,|b|) \Bigm| \lcm(2,6) = 6,
  \], so \(|(a,b)|\neq 4\). We conclude that \(\Z_4\times\Z_6/\<(2,3)\> \cong \Z_{12}\).
\end{example}

\section{Center Subgroups}

\begin{definition}
  The \emph{center} of a group \(G\) is the subgroup \[
    Z(G) = \{g\in G \mid gx = xg \text{ for all } x\in G\}.
  \]
\end{definition}

Clearly, \(G\) is abelian if and only if \(Z(G) = G\).

\begin{proposition}
  Let \(G\) be a group, then the center \(Z(G)\) is a normal subgroup of \(G\).
\end{proposition}
\begin{proof}
  \begin{itemize}
    \item Clealy \(e\in Z(G)\). Let \(g_1,g_2\in Z(G)\), then \(g_1x = xg_1\) and \(g_2x = xg_2\) for all \(x\in G\). Thus, \[
      (g_1g_2)x = g_1(g_2x) = g_1(xg_2) = (g_1x)g_2 = (xg_1)g_2 = x(g_1g_2),
    \] so \(g_1g_2\in Z(G)\), showing \(Z(G) < G\).
    \item Let \(g\in Z(G)\), \(a\in G\), then \(gx = xg\) for all \(x\in G\). Thus, \begin{align*}
      (aga^{-1})x = a(ga^{-1})x = a(a^{-1}g)x &= gx \\
      &= xg = x(ga)a^{-1} = x(ag)a^{-1} = x(aga^{-1}),
    \end{align*}
    \item so \(aga^{-1}\in Z(G)\), showing \(Z(G)\lhd G\). \qedhere
  \end{itemize}
\end{proof}

\newpage
\begin{lemma}
  If \(G/Z(G)\) is cyclic, then \(G\) is abelian.
\end{lemma}
\begin{proof}
  Let \(G/Z(G) = \<aZ(G)\>\) for some \(a\in G\). Let \(x,y\in G\), then \[
    xZ(G) = (aZ(G))^m = a^mZ(G), \quad
    yZ(G) = (aZ(G))^n = a^nZ(G)
  \] for some \(m,n\geq 0\), then \(x = a^m g\) and \(y = a^n g'\) for some \(g,g'\in Z(G)\). Thus \[
    xy = (a^m g)(a^n g') = a^ma^ngg' = a^na^mg'g
    = (a^ng')(a^mg) = yx. \qedhere
  \]
\end{proof}

\begin{example}
  For \(S_n\), \(n\neq 2\), we have \(Z(S_n) = \{\id\}\). On the other hand, \(Z(S_2) = S_2\).
\end{example}
\begin{proof}
  Clearly \(Z(S_n) = S_n\) for \(n\leq 2\) since they are abelian. Now suppose \(n\geq 3\).
  Fix \(\sgm\in Z(S_n)\), let \(i,j,k\in\{1,\ldots,n\}\) be distinct, note by Proposition~\ref{prop:S_n_transp_conjg} that \begin{align*}
    (i\ j) = \sgm(i\ j)\sgm^{-1} = (\sgm(i)\ \sgm(j)), \qquad
    (j\ k) = \sgm(j\ k)\sgm^{-1} = (\sgm(j)\ \sgm(k)),
  \end{align*}
  which gives \(\{\sgm(i),\sgm(j)\} = \{i,j\}\) and \(\{\sgm(j),\sgm(k)\} = \{j,k\}\), forcing \(\sgm(j) = j\). Then \(\sgm = \id\).
\end{proof}

\begin{example}
  For \(GL_n(\R)\), we have \(Z(GL_n(\R)) = \spn_\R^\times\{I_n\}\).
\end{example}
\begin{proof}
  Fix \(A\in Z(GL_n(\R))\), put \(E_{ij} = (\delta_{ij})_{ij}\), then \(E_{ij}^2 = \O_n\), so \[
    (1 - E_{ij})(1 + E_{ij}) = 1 - E_{ij}^2 = I_n,
  \] which gives \(1 + E_{ij}\in GL_n(\R)\). Now \[
    A(1 + E_{ij}) = (1 + E_{ij})A \implies AE_{ij} = E_{ij}A.
  \] Here \[
    (AE_{ij})_{kl} = \sum_{m=1}^n A_{km}(E_{ij})_{ml} = A_{ki}\delta_{jl}, \quad
    (E_{ij}A)_{kl} = \sum_{m=1}^n (E_{ij})_{km}A_{ml} = \delta_{ki}A_{jl},
  \] therefore \begin{itemize}
    \item put \((k,l) = (i,i)\) to see \(A_{ij} = 0\) for \(j\neq i\);
    \item put \((k,l) = (i,j)\) to see \(A_{ii} = A_{jj}\) for all \(i,j\).
  \end{itemize} Then \(A\in\spn^\times_\R(I_n)\).
\end{proof}

\begin{example}
  \(Z(D_n) = \begin{cases}
    \{e\}, & n \text{ odd} \\
    \{e,r^{n/2}\}, & n \text{ even}
  \end{cases}\).
\end{example}
\begin{proof}
  Let \(g\in Z(D_n)\), but in \(D_n\), we only have \(g = r^k\) or \(g = r^ks\).
  \begin{itemize}
    \item If \(g = r^k\), put \(x = s\in D_n\), then \[
      gx = xg \implies r^ks = sr^k \implies r^ks = r^{-k}s \implies r^{2k} = e \implies k = \frac{n}{2}.
    \]
    \item If \(g = r^ks\), put \(x = r\in D_n\), then \[
      gx = xg \implies r^ksr = rr^ks \implies r^{k-1}s = r^{k+1}s \implies r^2 = e \implies n\leq 2. \qedhere
    \]
  \end{itemize}
\end{proof}


It is in general not true that if \(H < G\) then \(Z(H) \subseteq H\cap Z(G)\), because in \(H\) there may be less elements needed to commute with than in \(G\).

\newpage
\begin{definition}
  The \emph{commutator} of two elements \(a,b\in G\) is the element \[
    [a,b] = aba^{-1}b^{-1}.
  \]
  The \emph{commutator subgroup} of a group \(G\) is the subgroup \[
    [G,G] = \<[a,b] \mid a,b\in G\>.
  \]
\end{definition}
Clearly, \(G\) is abelian if and only if \([G,G] = \{e\}\).

\begin{definition}
  The \emph{abelianization} of a group \(G\) is the quotient group \[
    G^{\ab} = G/[G,G].
  \]
\end{definition}

\begin{proposition}
  The commutator subgroup \([G,G]\) is the smallest normal subgroup of \(G\) such that \(G/[G,G]\) is abelian.
\end{proposition}

\begin{example}
  \([S_n,S_n] = A_n\) and \(S_n^\ab \cong \Z_2\).
\end{example}
\begin{proof}
  % Define \[
  %   \sgn: S_n \longrightarrow \Z_2, \quad
  %   \sgn(\tau) = \begin{cases}
  %     0, & \tau \text{ is even} \\
  %     1, & \tau \text{ is odd}
  %   \end{cases},
  % \] then we see \(\ker\sgn = A_n\) and \(\im\sgn = \Z_2\), so by the first isomorphism theorem, we have \[
  %   S_n/A_n \cong \Z_2,
  % \] which is abelian, so \([S_n,S_n]\leq A_n\). 
  Consider \[
    [(a\ b),(b\ c)] = (a\ b)(b\ c)(b\ a)(c\ b) = (a\ c\ b),
  \] then \([S_n,S_n]\) contains all \(3\)-cycles, which generate \(A_n\), so \([S_n,S_n]\geq A_n\). Note that since \([S_n:A_n] = 2\), we have \[
    S_n / A_n \cong \Z_2,
  \] which is abelian, so \([S_n,S_n]\leq A_n\). We conclude that \([S_n,S_n] = A_n\).
\end{proof}

\begin{example}
  \([D_n,D_n] = \<r^2\>\) and \[
    D_n^\ab \cong \begin{cases}
      \Z_2\times\Z_2, & n \text{ even} \\
      \Z_2, & n \text{ odd}
    \end{cases}.
  \]
\end{example}
\begin{proof}
  Consider \[
    [r,s] = rsr^{-1}s^{-1} = rrs^{-1}s^{-1} = r^2,
  \] then \([D_n,D_n]\geq \<r^2\>\).
  Note that \begin{align*}
    D_n/\<r^2\> = \<r^n = r^2 = s^2 = e, sr = r^{-1}s\>
  \end{align*} is abelian since every element has order at most \(2\). Then \([D_n,D_n] \leq \<r^2\>\).
  Furthermore, we compute \begin{align*}
    D_n^\ab = D_n/\<r^2\> \cong \begin{cases}
      \<r,s \mid r^2 = s^2 = e\> \cong \Z_2\times\Z_2, & n \text{ even} \\
      \<r,s \mid r\phantom{^2} = s^2 = e\> \cong \Z_2, & n \text{ odd}
    \end{cases}.
  \end{align*}
\end{proof}

\newpage

\begin{example}
  Let \[
    G = \Bigg\{
      \begin{pmatrix}
        a & b \\
        0 & c
      \end{pmatrix} \in \GL_2(\R) \Biggm| a, c \in \R^\times, b \in \R
    \Bigg\},
    \qquad
    N = \Bigg\{
      \begin{pmatrix}
        1 & b \\
        0 & 1
      \end{pmatrix} \in G \Biggm| b \in \R
    \Bigg\}.
  \]
  Then \([G,G] = N\).
\end{example}
\begin{proof}
  Consider \begin{align*}
    \Bigg[
      \begin{pmatrix}
        a & 0 \\
        0 & c
      \end{pmatrix},
      \begin{pmatrix}
        1 & b \\
        0 & 1
      \end{pmatrix}  
    \Bigg]
    &= \begin{pmatrix}
      a & 0 \\
      0 & c
    \end{pmatrix}
    \begin{pmatrix}
      1 & b \\
      0 & 1
    \end{pmatrix}
    \begin{pmatrix}
      1/a & 0 \\
      0 & 1/c
    \end{pmatrix}
    \begin{pmatrix}
      1 & -b \\
      0 & 1
    \end{pmatrix} \\
    &= \begin{pmatrix}
      1 & \frac{a}{c}b \\
      0 & 1
    \end{pmatrix}
    \begin{pmatrix}
      1 & -b \\
      0 & 1
    \end{pmatrix}
    = \begin{pmatrix}
      1 & (\frac{a}{c}-1)b \\
      0 & 1
    \end{pmatrix} \in N,
  \end{align*}
  so \([G,G]\geq N\).
  % Note that \(G\cong \R^\times\times\R\times\R^\times\) and \(N\cong \R\), we have \[
  %   \frac{G}{N} \cong \frac{\R^\times\times\R\times\R^\times}{\R} \cong \Bigg\{
  %     \begin{pmatrix}
  %       a & 0 \\
  %       0 & c
  %     \end{pmatrix} \in G \Biggm| a, c \in \R^\times
  %   \Bigg\} \cong \R^\times\times\R^\times,
  % \]
  Consider \[
    \f: G \longrightarrow \R^\times\times\R^\times,
    \quad
    \begin{pmatrix}
      a & b \\
      0 & c
    \end{pmatrix} \longmapsto (a-1,c-1),
  \] then \(\ker\f = N\) and \(\im\f = \R^\times\times\R^\times\), so \(G/N\cong\R^\times\times\R^\times\),
  which is abelian, so \([G,G]\leq N\).
\end{proof}


\section{Series of Groups}
Let \(G\) be a group.
\begin{definition}
  A \emph{subnormal series} of \(G\) is a finite chain of subgroups \[
    \{e\} = H_0 < H_1 < \cdots < H_n = G
  \] such that \(H_i \lhd H_{i+1}\) for all \(i\). We call \(n\) the \emph{length} of the series.
\end{definition}

\begin{definition}
  A subnormal series of \(G\) is called a \emph{normal series} if \(H_i \lhd G\) for all \(i\).
\end{definition}

\begin{example}
  \(\{0\} < 8\Z < 4\Z < \Z\).
\end{example}

\begin{example}
  \(\{\id\} < \<(1\ 2)(3\ 4)\> < V_4 < D_4\), where \[
    V_4 = \{\id, (1\ 2)(3\ 4), (1\ 3)(2\ 4), (1\ 4)(2\ 3)\}\cong \Z_2\times\Z_2
  \] is called the \emph{Klein four-group}.
\end{example}

\begin{definition}
  Given two subnormal series \(\{H_i\}_i\) and \(\{K_j\}_j\) of a group \(G\). We say \(\{K_j\}_j\) is a \emph{refinement} of \(\{H_i\}_i\) if \[
    \{H_0, H_1, \ldots, H_n\} \subset \{K_0, K_1, \ldots, K_m\}
  \]
\end{definition}

\begin{definition}
  A subnormal (resp. normal) series \(\{H_i\}_i\) of \(G\) is called a \emph{composition} (resp \emph{principal}) \emph{series} if it has no proper refinement.
\end{definition}
Equivalently, a composition series is a subnormal series \(\{H_i\}_i\) such that \(H_{i+1}/H_i\) is simple for all \(i\).

\begin{example}
  The principal series of \(S_n\) when \begin{itemize}
    \item \(n \geq 5\) or \(n = 3\), is given by \[
      \{e\} < A_n < S_n.
    \]
    \item \(n = 4\), is given by \[
    \{e\} < \<(1\ 2)(3\ 4)\> < V_4 < A_4 < S_4.
  \]
  \end{itemize}
\end{example}

Clearly, every finite group has a composition series. Given a composition series \[
  \{e\} = H_0 < H_1 < \cdots < H_n = G,
\] we obtain a sequence of short exact sequences \[
  \begin{tikzcd}[row sep=tiny]
    1 \ar{r} & H_1 \ar[hook]{r} & H_2 \ar[->>]{r} & H_2/H_1 \ar{r} & 1 \\
    1 \ar{r} & H_2 \ar[hook]{r} & H_3 \ar[->>]{r} & H_3/H_2 \ar{r} & 1 \\
    & & \vdots \\
    1 \ar{r} & H_{n-1} \ar[hook]{r} & H_n \ar[->>]{r} & G/H_n \ar{r} & 1
  \end{tikzcd}
\]

Here \(1\) denotes the trivial group. Given groups \(G, H, Q\), a sequence of maps \[
  \begin{tikzcd}
    1 \ar{r} & H \ar[hook]{r}{\iota} & G \ar[->>]{r}{\pi} & Q \ar{r} & 1
  \end{tikzcd}
\] is called \emph{short exact} if \begin{enumerate}
  \item \(\iota\) is injective;
  \item \(\pi\) is surjective;
  \item \(\ker\pi = \im\iota\).
\end{enumerate} In this case, we say \(G\) is an \emph{extension} of \(H\) by \(Q\).

\begin{theorem}[Jordan-Hölder theorem]
  Let \(G\) be a finite group. Suppose \begin{align*}
    \{e\} = H_0 < H_1 < \cdots < H_n = G, \\
    \{e\} = K_0 < K_1 < \cdots < K_m = G
  \end{align*} are two composition series for \(G\), then \(n = m\) and \[
    H_{i+1}/H_i \cong K_{\sgm(i)+1}/K_{\sgm(i)}
  \] for some \(\sgm\in S_n\).
\end{theorem}

\section{Solvable Groups}

\begin{definition}
  A group \(G\) is called \emph{solvable} if it has a subnormal series \[
    \{e\} = H_0 < H_1 < \cdots < H_n = G
  \] such that each factor \(H_{i+1}/H_i\) is abelian. In this case, we call such series a \emph{solvable series}.
\end{definition}
In other words, a solvable group is a group built from successive extensions of abelian groups.

\begin{example}
  Every abelian group is solvable, since we can take the subnormal series \(\{e\} < G\).
\end{example}

\begin{example}
  \(S_3\) and \(S_4\) and their subgroups are solvable, since we have the subnormal series \[
    \{e\} < A_3 < S_3, \qquad
    \{e\} < V_4 < A_4 < S_4.
  \]
\end{example}

\begin{example}
  \(S_n\) is not solvable for \(n\geq 5\), since \(A_n\) is simple and nonabelian.
\end{example}

\begin{proposition}\
  \begin{enumerate}
    \item Every subgroup and quotient of a solvable group is solvable;
    \item If \(N\lhd G\) and both \(N\) and \(G/N\) are solvable, then \(G\) is solvable.
  \end{enumerate}
\end{proposition}
\begin{proof}
  \begin{enumerate}[label=(\alph*)]
    \item Suppose that \(G\) has a solvable series \[
      \{e\} = K_0 < K_1 < \cdots < K_n = G.
    \] Let \(H < G\) be a subgroup. Then consider \[
      \f:H\cap K_{i+1} \longrightarrow K_{i+1}/K_i, \quad
      h \longmapsto hK_i.
    \] Now \(\ker\f = H\cap K_i\) and \(\im\f = H\cap K_{i+1}/K_i\) and \[
      \frac{H\cap K_{i+1}}{H\cap K_i} \cong \frac{H\cap K_{i+1}}{K_i} < \frac{K_{i+1}}{K_i},
    \] so we have a solvable series \[
      \{e\} = H\cap K_0 < H\cap K_1 < \cdots < H\cap K_n = H.
    \]
  \end{enumerate}
\end{proof}

\begin{example}
  A group \(G\) of order \(38\) is solvable.
\end{example}
\begin{proof}
  By Cauchy's theorem, \(G\) admits a subgroup \(H \cong \Z_{19}\). Since \([G:H] = \frac{38}{19} = 2\), we have \(H\) is normal. Therefore, we have a normal series \[
    \{e\} < H < G,
  \] where \(H\) and \(G/H\) are abelian, so \(G\) is solvable.
\end{proof}
  
\section{Derived Series}

The commutator subgroup \([G,G]\) of \(G\) is also called the \emph{first derived subgroup} of \(G\), denoted by \(G^{(1)}\). For \(k\geq 1\), we define the \(k\)-th derived subgroup of \(G\) recursively by \[
  G^{(k)} := [G^{(k-1)}, G^{(k-1)}].
\]

\begin{definition}
  We call the series \[
    G > G^{(1)} > G^{(2)} > \cdots
  \] the \emph{derived series} of \(G\).
\end{definition}

\begin{example}
  \(S_n > A_n > A_n > \cdots\).
\end{example}

\begin{proposition}
  A group \(G\) is solvable if and only if \(G^{(k)} = \{e\}\) for some \(k\).
\end{proposition}
\begin{proof}
  For \((\Leftarrow)\), consider the dervied series \[
    G > G^{(1)} > \cdots > G^{(k)} = \{e\},
  \] then each \(G^{(i-1)}/G^{(i)} = (G^{(i-1)})^\ab\), which is abelian, so \(G\) is solvable.

  For \((\Rightarrow)\), if we have a solvable series \[
    \{e\} = H_0 < H_1 < \cdots < H_n = G,
  \] then since \(G/H_{n-1}\) is abelian, we have \(G^{(1)}\leq H_{n-1}\).
  Now \(G^{(1)}H_{n-2} < H_{n-1}\), then by the second isomorphism theorem, we have \[
    \frac{G^{(1)}}{G^{(1)}\cap H_{n-2}} \cong \frac{G^{(1)}H_{n-2}}{H_{n-2}} < \frac{H_{n-1}}{H_{n-2}},
  \] which is abelian, so \(G^{(2)} < G^{(1)}\cap H_{n-2} < H_{n-2}\), repeating this argument shows that \(G^{(i)} < H_{n-i}\) for all \(i\). Eventually, \(G^{(k)} = \{e\}\) for some \(k\leq n\).
\end{proof}

\begin{example}
  The group \(U_2(\R) = \Bigg\{
      \begin{pmatrix}
        a & b \\
        0 & c
      \end{pmatrix} \in \GL_2(\R) \Biggm| a, c \in \R^\times, b \in \R
    \Bigg\}\) is solvable.
\end{example}
\begin{proof}
  Recall that \(G^{(1)} = \Bigg\{
      \begin{pmatrix}
        1 & b \\
        0 & 1
      \end{pmatrix} \in G \Biggm| b \in \R
    \Bigg\}\), which is abelian, so \(G^{(2)} = \{I_2\}\).
    Therefore, \(G\) is solvable.
\end{proof}